\subsection{The unresolved relative-period bound}
\hypertarget{tag-001D}{}\label{ari:C.7.4}

The open lower-bound question is the Relative Semi-abelian Conjecture
\cite[Conjecture~4.7, arXiv v2, p.~18]{bertolin2022} which asks whether algebraic lifts in the indicated semiabelian
varieties satisfy the bound
\[
 \operatorname{trdeg}_{\overline{\mathbb Q}(\Omega_G)}
 \overline{\mathbb Q}\bigl(\Omega_G,R,\bar R,
                   \widetilde{\log}_G(R,\bar R)\bigr)
 \geq 2\dim B_Q+\dim Z(1).
\]
The full matrix $\Omega_G$ includes the second-kind periods.
This lower bound remains open for the compact diagonal
configuration at issue here. Its relation to Grothendieck's
period conjecture for algebraic $1$-motives is given in
\cite[Proposition~4.12, arXiv v2, pp.~19--20]{bertolin2022}.

For the algebraic motives of \cref{ari:C.7.3}, projection
from the point-extension variety to the extension-coordinate
variety is an isogeny: the point coordinates agree with the
extension data under the stated polarizations.
Equivalently, $(P,\bar P)=(P,0)+(0,\bar P)$ in
$E\times\bar E$. Formula~(4.4), p.~19 of the same source
therefore gives $\dim B_Q=0$.
The remaining kernel of connected motivic groups is the
weight-$-2$ vector group $\operatorname{Lie}Z$.
After fixing the base periods its fibre is an affine torsor
of dimension $\dim Z$.

In positive Tate coordinates the central action changes
$(v,\bar v)$ by $2\pi i(c_1,c_2)$.
For an isogeny $\phi:\bar E\to E$ of degree $d$,
$\phi^{\mathsf t}J_E\phi=dJ_{\bar E}$.
Thus the dependent bracket line of \cref{ari:C.7.3}
has coordinates $(1,aa^\dagger/d)$ with positive second
coordinate; in the independent case the bracket spans both
central coordinates.
The functional $(c_1,c_2)\mapsto c_1+c_2$ is nonzero
on this central group, including when there is a further
quotient Tate direction.
Accordingly, the compactness equation
\eqref{ari:eq:cp-relation} cuts a proper affine hyperplane.
The unconditional period-field calculation is
\cref{I:compact-period-relation}; the proposed lower bound
needed to exclude an algebraic compact lift is the open step.

